\begin{minipage}{.3\textwidth}
  \flushleft
  \center{\includegraphics[scale=.09]{unam.pdf}}

  \vspace{20pt}

  \center{
    \rule{.5pt}{.6\textheight}
    \hspace{7pt}
    \rule{2pt}{.6\textheight}
    \hspace{7pt}
    \rule{.5pt}{.6\textheight}
  } \\

  \center{\includegraphics[scale=.22]{ciencias.pdf}}
\end{minipage}
\begin{minipage}{.7\textwidth}
\flushright

\center{

  \center{
    \LARGE{U}\large{NIVERSIDAD} \LARGE{N}\large{ACIONAL}
    \LARGE{A}\large{UTÓNOMA} \\[10pt]
    \large{DE}
    \LARGE{M}\large{ÉXICO}
  } \\
  \rule{\textwidth}{2pt}
  \\
  \hrulefill\\[1cm]

  \LARGE{F}\large{ACULTAD DE } \LARGE{C}\large{IENCIAS}\\[2cm]

  \large{
  Análisis del algoritmo \textit{Knuth-Morris-Pratt} puramente funcional
  }\\[1.6cm]

  \huge{
T \hspace{1cm} E \hspace{1cm} S \hspace{1cm} I \hspace{1cm} S  }\\[1cm]

  \large{QUE PARA OBTENER EL TÍTULO DE:}\\[1cm]

  \large{
Licenciado en Ciencias de la Computación  }\\[1cm]

  \large{PRESENTA:}\\[1cm]

  \large{
Ángel Iván Gladín García  }\\[1cm]

  \large{
TUTORA  }\\[1cm]

  \large{
  Dra. Lourdes Del Carmen Gonzales Huesca}\\
  \large{
    Ciudad Universitaria, Cd. Mx., 2020
  }
}

\end{minipage}
